\documentclass[11pt]{article}
\usepackage[margin=1in]{geometry}
\usepackage{amsmath,amssymb,booktabs,xcolor}
\usepackage[hidelinks]{hyperref}
\newcommand{\E}{\mathbb E}
\newcommand{\todo}[1]{\textcolor{blue}{[TODO: #1]}}
\title{Ordered Non-Reset Power-Grid sBOED: EIG, MSC and MOCU}
\author{}
\date{}
\begin{document}
\maketitle
This specification matches the active online IEEE9 experiment. Probe duration
and terminal control are continuous but distinct decisions. The physical system
is initialized only once per episode. The reset design catalogs and banks belong
to the historical protocol and are not used for this specification.
\section{Problem Formulation}
\label{sec:problem}
\subsection{Power-Grid Dynamics and Latent Parameter Space}
Following Peng et al.~\cite{peng2024probing}, we model inertia and fast
frequency response jointly. Let $\Delta\omega(s)$ be angular-speed deviation
from an equilibrium with synchronous speed $\omega_{\mathrm s}=2\pi f_0$.
On a common power base $S_{\mathrm{base}}$, the incremental power balance is
\begin{equation}
M\frac{d\Delta\omega}{ds}
=\Delta P_m-\Delta P_e-D\Delta\omega,
\qquad M=\frac{2HS_n}{\omega_{\mathrm s}S_{\mathrm{base}}}.
\label{eq:inertia-balance}
\end{equation}
Here $H$ is the inertia constant, $S_n$ the rating associated with $H$,
$\Delta P_m$ and $\Delta P_e$ the input and electrical-output power changes,
and $D$ the damping coefficient. Power changes are in per unit.
The paper represents fast frequency support by a droop relation,
\begin{equation}
\Delta P_m=-K\Delta f=-\frac{K}{2\pi}\Delta\omega,
\qquad \Delta f=\frac{\Delta\omega}{2\pi},
\label{eq:frequency-support}
\end{equation}
where $K$ is expressed here in per-unit power per Hz. Substitution gives
\begin{equation}
M\frac{d\Delta\omega}{ds}
=-\Delta P_e-\left(\frac{K}{2\pi}+D\right)\Delta\omega.
\label{eq:aggregate_model}
\end{equation}
Thus $M$ governs the rate of frequency change, whereas $K$ governs the
restoring power supplied in response to a frequency deviation. We retain
known damping $D$; Peng et al. neglect it in their estimation model.

To represent probe location, we extend this aggregate relation to a reduced
network. Let $\mathcal N$ be the physical-bus set and
$\mathcal M\subseteq\mathcal N$ the dynamic-machine buses, with
$m=|\mathcal M|$. For $i\in\mathcal M$, write $\omega_i$ for angular-speed
deviation and $\delta_i$ for rotor angle in the synchronous frame:
\begin{align}
\dot\delta_i(s)&=\omega_i(s),\\
M_i\dot\omega_i(s)&=-\Delta P_{e,i}(\boldsymbol\delta(s))
-\left(\frac{K_i}{2\pi}+D_i\right)\omega_i(s)+v_i(s).
\label{eq:swing-frequency}
\end{align}
The initial state is $(\boldsymbol\delta^0,\mathbf0)$, and the reduced
lossless-network power change is
\begin{equation}
\Delta P_{e,i}(\boldsymbol\delta)=\sum_{j\in\mathcal M}B_{ij}
\left[\sin(\delta_i-\delta_j)-\sin(\delta_i^0-\delta_j^0)\right],
\label{eq:network-power}
\end{equation}
with known coupling coefficients $B_{ij}$. A physical-bus injection is
mapped to $\mathbf v$ by the known vector $\mathbf a(b)$; the reduction is
given in Eq.~\eqref{eq:kron-map}. This machine-wise model is our extension
of the paper's aggregate response, with $K_i$ representing effective fast
frequency support at each retained machine bus.

The latent vector contains the inertia and response coefficients:
\begin{equation}
\boldsymbol\theta=(M_1,\ldots,M_m,K_1,\ldots,K_m)^\top\in\Theta,
\qquad\boldsymbol\theta\sim p_0,
\label{eq:theta-prior}
\end{equation}
where indices $1,\ldots,m$ follow the retained-machine ordering and
\begin{equation}
\Theta=\prod_{i=1}^{m}[M_i^-,M_i^+]
\times\prod_{i=1}^{m}[K_i^-,K_i^+],
\qquad M_i^->0,\quad K_i^-\ge0.
\label{eq:latent-space}
\end{equation}
The topology, coupling, damping, and equilibrium are known. Algebraic buses
enter the coupling and injection map but have no independent $M_i$ or $K_i$.
With bounded piecewise-continuous inputs, these assumptions define the
frequency trajectories used by the observation and terminal-control models.

\subsection{Ordered Non-Reset Probes and Observations}
The controllable probe variable is its duration, $\xi_t=\tau_t$.
Injection bus $b_{\mathrm p}$ and amplitude $A_{\mathrm p}$ are fixed.
A stage occupies a recording window of length $W$, independently of the
injection duration. With $s_t=(t-1)W$ and local time $\ell=s-s_t$, the
stage input is
\begin{equation}
\mathbf v_{\mathrm p,t}(\ell)
=\mathbf a(b_{\mathrm p})\frac{A_{\mathrm p}}{2}
\left[1-\cos\left(\frac{2\pi\ell}{\tau_t}\right)\right]
\mathbf1_{[0,\tau_t]}(\ell),\quad 0\le\ell\le W.
\label{eq:probe-input}
\end{equation}
Let $\mathbf x=(\boldsymbol\delta,\boldsymbol\omega)$ and let
$\Phi_\theta(\mathbf x,\tau;\ell)$ denote propagation of the swing equations
under this input. The state transition is
\begin{equation}
\mathbf x_t(\theta)=\Phi_\theta(\mathbf x_{t-1}(\theta),\tau_t;W),
\qquad \mathbf x_0=(\boldsymbol\delta^0,\mathbf0).
\label{eq:carry}
\end{equation}
Only the beginning of a new episode uses $\mathbf x_0$. Later stages use the
preceding terminal state. The latent $M_i,K_i$ do not change within an episode.
Consequently, $(\tau_1,\tau_2)$ and $(\tau_2,\tau_1)$ are different input
histories, even when their unordered duration sets coincide.

Write $h_t=((\tau_1,\mathbf y_1),\ldots,(\tau_t,\mathbf y_t))$ for the
ordered available history. The design space at stage $t$ is
\begin{equation}
\Xi(h_{t-1})=\{\tau\in[\tau_{\min},\tau_{\max}]:
|\tau-\tau_k|\ge\varepsilon_\tau\text{ for every }k<t\},
\label{eq:design-space}
\end{equation}
where $0<\tau_{\min}<\tau_{\max}\le W$. This is a continuous set with
previously used neighborhoods excluded, not a catalog of bus--duration pairs.

At known local sample times $0<r_1<\cdots<r_{N_{\mathrm{obs}}}\le W$, a frequency sensor
at physical bus $b_{\mathrm y}$ produces
\begin{equation}
\mathbf y_t=\mathbf g(\theta,\mathbf x_{t-1}(\theta),\tau_t)
+\boldsymbol\epsilon_t,
\qquad \boldsymbol\epsilon_t\sim\mathcal N(\mathbf0,\sigma^2\mathbf I_{N_{\mathrm{obs}}}),
\label{eq:observation}
\end{equation}
where $j=1,\ldots,N_{\mathrm{obs}}$ and $g_j=\Delta f_{b_{\mathrm y}}(s_t+r_j)$. Independent sensor noise is
added to observations, not to the propagated physical state. The last sample
is available before the next design decision. The benchmark assumes zero
computation/communication delay; a delayed implementation would also have to
propagate the state during that delay.

\subsection{History-Conditioned Bayesian Update and EIG}
The posterior is
\begin{equation}
p_t(\theta)=
\frac{p_{t-1}(\theta)
 p(\mathbf y_t\mid\theta,h_{t-1},\tau_t)}
{\int p_{t-1}(\vartheta)
 p(\mathbf y_t\mid\vartheta,h_{t-1},\tau_t)\,d\vartheta},
\label{eq:posterior}
\end{equation}
with the Gaussian likelihood from Eq.~\eqref{eq:observation}. Its mean must
be simulated from the candidate model's own carried state. An equilibrium
response indexed only by $(\theta,\tau_t)$ is insufficient after the first probe.

For a deterministic policy $\tau_t=\pi_\phi(h_{t-1})$, the information objective is
\begin{equation}
J_{\mathrm{EIG}}(\pi_\phi)=
\E_{\theta\sim p_0,\,H_T\mid\theta,\pi_\phi}
\left[\log\frac{p(\theta\mid H_T)}{p_0(\theta)}\right].
\label{eq:eig}
\end{equation}
The implementation uses the sequential prior-contrastive lower bound
(sPCE) of DAD~\cite{foster2021dad}, with history-conditioned likelihoods:
\begin{align}
\ell_t^{(n)}&=\sum_{k=1}^t
\log p(\mathbf y_k\mid\theta^{(n)},h_{k-1},\tau_k),\\
\widehat I_t&=\ell_t^{(0)}-
\log\left[\frac{1}{L+1}\sum_{n=0}^{L}\exp(\ell_t^{(n)})\right].
\label{eq:spce}
\end{align}
Here $\theta^{(0)}$ generates the observations and the other models are
independent prior draws. The reported expected bound depends on $L$ and
has ceiling $\log(L+1)$. These scores are not pooled with historical
bank-entropy estimates.

\subsection{Continuous Terminal Control and Physical Safety}
The final engineering decision is an injection magnitude
\begin{equation}
u_{\mathrm{ctrl}}\in\mathcal U=[u_{\min},u_{\max}],
\label{eq:control-space}
\end{equation}
which is continuous and distinct from the probe duration in
Eq.~\eqref{eq:design-space}. At global time $TW$, the probe input ends and
the declared contingency and supplementary step control begin. With local
control time $r$, the input over the control-assessment window $[0,H_c]$ is
\begin{equation}
\mathbf v_{\mathrm c}(r;u)=\mathbf p_{\mathrm c}+\mathbf a(b_{\mathrm c})u,
\qquad \mathbf x_{\mathrm c}(0;\theta,h_T)=\mathbf x_T(\theta).
\label{eq:control-input}
\end{equation}
The disturbance, injection location, window and units are fixed before design
optimization. Thus the connection from probing to action is explicit:
observations update the posterior over $M/K$, each candidate $M/K$ determines
a carried terminal state, and the controller evaluates safety from those states.

For the monitored dynamic buses $\mathcal M$, define
\begin{align}
f_{\min}(\theta,h,u)&=\min_{i\in\mathcal M,\,0\le r\le H_c}
\Delta f_i(r;\theta,\mathbf x_T(\theta),u),\\
R_{\max}(\theta,h,u)&=\max_{i\in\mathcal M,\,0\le r\le H_c}
\left|\frac{d\Delta f_i(r;\theta,\mathbf x_T(\theta),u)}{dr}\right|,\\
S(\theta,h,u)&=\mathbf1\{f_{\min}(\theta,h,u)\ge\underline f,
\ R_{\max}(\theta,h,u)\le\overline R\}.
\label{eq:safety}
\end{align}
The initial frequency is included in this window. These are declared
frequency-security criteria, not a claim of general transient stability or
certified safety throughout the preceding probes.
The model-specific minimum safe control is
\begin{equation}
U(\theta,h)=\inf\{u\in\mathcal U:S(\theta,h,u)=1\},
\label{eq:oracle}
\end{equation}
with $U=+\infty$ when the admissible safe set is empty.
The current scalar solver assumes that a nonempty safe set is
$[U(\theta,h),u_{\max}]$. It checks sampled monotonicity and rejects a detected
violation; those numerical checks do not prove the assumption globally.

\subsection{Minimum Safe Control as a Design Objective}
For a specified posterior coverage $q\in(0,1)$, select
\begin{equation}
u_{\mathrm{MSC}}(h;q)=\inf\left\{u\in\mathcal U:
\int S(\theta,h,u)p(\theta\mid h)\,d\theta\ge q\right\}.
\label{eq:msc-controller}
\end{equation}
If this set is empty, the decision is infeasible. Under the scalar safety
assumption, Eq.~\eqref{eq:msc-controller} is the posterior $q$ quantile of
$U(\theta,h)$. Coverage expresses a conditional model-based requirement;
it is not an observed safety rate or a value established by an engineering standard.

We use \emph{MSC} as a descriptive abbreviation for minimum safe control,
not as a claim of standardized terminology. Its BOED target is
\begin{equation}
\min_\pi J_{\mathrm{MSC}}(\pi),\qquad
J_{\mathrm{MSC}}(\pi)=\E_{H_T\mid\pi}[u_{\mathrm{MSC}}(H_T;q)].
\label{eq:msc-objective}
\end{equation}
The physical requirement $U(\theta^\star,h)$ evaluated with true parameters
is the oracle; $u_{\mathrm{MSC}}(h;q)$ is the implementable decision using
noisy observations. They are not interchangeable. Lower expected MSC values
smaller selected control at the declared posterior coverage, without an
under-control penalty in the design score.

\subsection{Finite-Loss Control-Oriented MOCU}
% BEGIN SHARED FINITE-LOSS MOCU
MOCU compares the loss of a decision made under uncertainty with that of a
model-informed decision~\cite{yoon2013mocu,boluki2018generalized}.
Our particular control-sizing loss, for finite feasible $U(\theta,h)$, is
\begin{equation}
C_q(u,\theta;h)=u+\frac{[U(\theta,h)-u]_+}{1-q}.
\label{eq:operational_cost}
\end{equation}
This is a chosen surplus/shortfall surrogate in control-magnitude units.
It is not an engineering-calibrated cost of violating the frequency or RoCoF
limits. The cited MOCU framework supplies the decision-theoretic construction,
not this particular loss or an engineering value of $q$.

For fixed $h$, the model-informed optimum is $u=U(\theta,h)$ with loss
$U(\theta,h)$. A posterior Bayes action satisfies
\begin{equation}
u_q(h)\in\arg\min_{u\in\mathcal U}
\E_{\theta\mid h}[C_q(u,\theta;h)].
\label{eq:bayes-action}
\end{equation}
Writing $F_h$ for the posterior CDF of $U(\theta,h)$, its optimality condition is
\begin{equation}
F_h(u_q^-)\le q\le F_h(u_q).
\label{eq:quantile_optimality}
\end{equation}
Choosing the smallest such action gives the same controller as
Eq.~\eqref{eq:msc-controller} under the stated threshold assumption. The
\emph{design objectives} nevertheless differ: the posterior MOCU is
\begin{equation}
\mathcal M(h;q)=\E_{\theta\mid h}
\left[u_q(h)+\frac{[U(\theta,h)-u_q(h)]_+}{1-q}-U(\theta,h)\right],
\label{eq:mocu}
\end{equation}
and MOCU-based design minimizes
\begin{equation}
J_{\mathrm{MOCU}}(\pi)=\E_{H_T\mid\pi}[\mathcal M(H_T;q)].
\label{eq:mocu-objective}
\end{equation}
MSC scores selected control; MOCU scores expected excess loss relative to
the model-specific requirement. MOCU is nonnegative under its assumptions
and vanishes when the required control is known exactly. A lower MOCU need
not imply a smaller selected control or a larger empirical safety rate.
Infeasible requirements make this finite-loss construction unsupported and
must not be hidden by clipping or resampling.

In the non-reset application, probing changes both the posterior and the
physical state entering control. Thus $U(\theta,h)$ itself depends on the
applied probe history. A fixed-application value-of-information monotonicity
argument cannot be transferred without qualification: increasing the horizon
need not reduce either MSC or MOCU, and any improvement is not automatically
attributable to information alone.
% END SHARED FINITE-LOSS MOCU

\subsection{Physical Safety Rate and Held-Out Evaluation}
The policy safety probability is
\begin{equation}
\mathrm{SR}(\pi;q)=\E_{\theta\sim p_0,\,H_T\mid\theta,\pi}
[S(\theta,H_T,u_q(H_T))].
\label{eq:safety-rate}
\end{equation}
With an exact posterior and the same correctly specified prior-predictive
model, conditional coverage in Eq.~\eqref{eq:msc-controller} implies
$\mathrm{SR}(\pi;q)\ge q$ by iterated expectation. A finite-particle
approximation, model mismatch or a finite test sample does not provide that
empirical guarantee.
For $N$ independent test systems and $R_i$ observation-noise replicas per system,
\begin{equation}
\widehat{\mathrm{SR}}=\frac1N\sum_{i=1}^N\frac1{R_i}
\sum_{r=1}^{R_i}S(\theta_i^\star,h_{i,r},u_q(h_{i,r})).
\label{eq:empirical-safety}
\end{equation}
Controls and posterior objectives use the same within-system averaging.
The held-out oracle is recomputed using $\theta_i^\star$ and that method's
actual terminal state. Different probe histories can therefore have different
oracle values even for the same true $M/K$. True parameters are used only
to generate/evaluate physical outcomes, never as controller or policy inputs.

\section{Design Methods and Numerical Implementation}
\begin{table}[t]
\caption{Method capabilities; these properties do not establish a performance advantage.}
\centering
\begin{tabular}{lccc}
\toprule
Method & \shortstack{Uses\\observations} & \shortstack{Horizon-aware\\objective} & \shortstack{Online policy\\refinement}\\
\midrule
Random & No & No & No\\
Fixed & No & Yes & No\\
Myopic & Yes & No & No\\
DAD & Yes & Yes & No\\
RL-sBOED & Yes & Yes & No\\
Step-DAD & Yes & Yes & Yes\\
\bottomrule
\end{tabular}
\end{table}

\subsection{Common Objective and Ordered History}
Define terminal utility $G_T$ as $\widehat I_T$ for EIG,
$-u_{\mathrm{MSC}}(h_T;q)$ for MSC, and $-\widehat{\mathcal M}(h_T;q)$ for
MOCU. Each objective defines a separate policy-training problem; an EIG
checkpoint is not reused as an MSC/MOCU policy. Inputs retain the stage
positions of all observed duration--frequency pairs. Exchangeable pooling
is inappropriate because Eq.~\eqref{eq:carry} depends on input order.
Every method uses the same admissible duration set and terminal controller.

DAD uses a deterministic map $\tau_t=\pi_\phi(h_{t-1})$, learned offline
by maximizing expected full-horizon utility. Its weights stay fixed during
an evaluated episode. This preserves the policy-based, non-myopic construction
of Foster et al.~\cite{foster2021dad}; the control objectives and numerical
training described below are explicit extensions of its EIG setting.

RL-sBOED uses a stochastic policy and a REDQ-style off-policy actor--critic,
with replay, a critic ensemble, random-subset target minimum, target-network
averaging and entropy regularization~\cite{blau2022rlsboed}. Let $G_t$ be
the corresponding utility assessed hypothetically at prefix $h_t$, with
$G_0=0$. The step rewards obey
\begin{equation}
r_t=G_t-G_{t-1},\qquad \gamma=1,
\qquad \sum_{t=1}^T r_t=G_T.
\label{eq:rl-rewards}
\end{equation}
For MSC/MOCU, a prefix control is a counterfactual used in the training
reward; physical control is applied only after stage $T$. The unregularized
return therefore targets the same terminal utility as DAD. REDQ additionally
uses its entropy regularizer during training; evaluation uses the actor's
mean transformed to the feasible duration set.

Step-DAD starts with its matching DAD policy. After observing a prefix, it
updates the posterior and refines a private policy by simulating the
\emph{remaining} horizon from posterior parameter/state particles. Past
inputs and observations remain fixed. The first probe uses the original
policy; later refinements retain previous within-episode updates. This
preserves the infer--refine structure of Hedman et al.~\cite{hedman2025stepdad}.

Fixed optimizes a deterministic ordered sequence offline, without access to
future observations. Its initialization is selected among short, central,
long and spread feasible sequences using validation data. Myopic searches
for the next duration maximizing expected next-stage utility under its
current posterior. Random samples uniformly over the lengths of the remaining
feasible intervals. Neither structural adaptivity nor horizon awareness
imposes an expected ordering of these methods' results.

\subsection{Online Posterior and Continuous Control Solver}
For posterior particles $\theta^{(n)}$, all states are propagated under the
actual ordered durations. Weights use
\begin{equation}
w_n^{(t)}\propto w_n^{(t-1)}
p(\mathbf y_t\mid\theta^{(n)},h_{t-1},\tau_t),\qquad \sum_nw_n^{(t)}=1.
\label{eq:particle-update}
\end{equation}
The controller's support is sampled independently of the system generating
an evaluation trajectory; that generating model is not deliberately inserted
into this posterior. Its role in the sPCE denominator is specific to
Eq.~\eqref{eq:spce} and must not be transferred to control inference.

For each carried state, a coarse numerical scan brackets the first detected
unsafe-to-safe transition, followed by bisection to approximate
$U_n=U(\theta^{(n)},h)$. The scan points are not admissible-action restrictions:
final controls take continuous values between them. A detected non-monotone
safe set stops the scalar solver. With normalized weights,
\begin{align}
\widehat F_h(v)&=\sum_n w_n\mathbf1\{\widehat U_n\le v\},\\
\widehat u_q(h)&=\inf\{v\in\mathcal U:\widehat F_h(v)\ge q\},\\
\widehat{\mathcal M}(h;q)&=\sum_nw_n
\left[\widehat u_q+\frac{(\widehat U_n-\widehat u_q)_+}{1-q}-\widehat U_n\right].
\label{eq:particle-objectives}
\end{align}
The selected action is checked again by direct safety simulation across the
posterior support. For MSC, an infeasible particle keeps its original
posterior mass with $\widehat U_n=+\infty$; a finite decision exists only if
coverage can still be met. Finite-loss MOCU instead requires finite feasible
requirements. No failed system is discarded to improve reported safety.

The full evaluation records ordered actions, noisy observations, posterior
weights and states, selected continuous control, true-system safety and a
separately computed continuous numerical oracle at the actual terminal state.
An equilibrium control bank indexed only by $M/K$ cannot replace this calculation.

\subsection{Gradient Treatment and Method Fidelity}
For EIG, deterministic DAD, Fixed and Step-DAD use pathwise differentiation
of the complete observation/state trajectory. Simulator state and duration
Jacobians are numerical central differences, with one-sided differences at
bounds. The feasible-duration map is differentiated locally within its
selected intervals. Interval switches and numerical derivatives require
sensitivity checks; no globally exact derivative is claimed.

Continuity of the physical action $u$ does not make a finite-particle empirical
quantile smooth in its posterior weights. The control-objective implementation
therefore retains deterministic design policies but uses antithetic numerical
policy-parameter perturbations~\cite{salimans2017es}. If $J(\phi)=\E[G_T]$,
$z\sim\mathcal N(0,I)$ and $d>0$ is a numerical optimization scale, the estimator is
\begin{equation}
\widehat{\nabla J_d}(\phi)=\frac1K\sum_{k=1}^K
\frac{\widehat J(\phi+d z_k)-\widehat J(\phi-d z_k)}{2d}\,z_k.
\label{eq:numerical-gradient}
\end{equation}
Each pair uses common model/noise draws and the same exact, numerically
refined terminal objective. This targets the Gaussian-smoothed parameter
objective $J_d(\phi)=\E_z[J(\phi+dz)]$, not an unbiased derivative of
unsmoothed $J$ at finite $d$. Validation and evaluation use the unperturbed
policy and the original controller. Scale/direction sensitivity is required
before performance claims. This is a declared optimizer adaptation, not
original DAD's pathwise estimator and not stochastic-action REINFORCE relabeled
as deterministic DAD. RL-sBOED retains its REDQ updates and does not require
terminal-controller derivatives.

\section{Experimental Protocol}
\subsection{Reduced-network construction}
Let $\mathcal A=\mathcal N\setminus\mathcal M$. For a connected lossless
network with fixed voltage magnitudes, let
$\mathbf L$ be the weighted susceptance Laplacian, ordered as
$(\mathcal M,\mathcal A)$. If $\mathcal A\ne\varnothing$, assume
$\mathbf L_{\mathcal A\mathcal A}$ is invertible and define
\begin{align}
\mathbf L_{\mathrm{red}}&=\mathbf L_{\mathcal M\mathcal M}
-\mathbf L_{\mathcal M\mathcal A}\mathbf L_{\mathcal A\mathcal A}^{-1}
\mathbf L_{\mathcal A\mathcal M},\\
\mathbf A_{\mathrm{inj}}&=
\left[\mathbf I_m\; -\mathbf L_{\mathcal M\mathcal A}
\mathbf L_{\mathcal A\mathcal A}^{-1}\right],\\
B_{ij}&=-(L_{\mathrm{red}})_{ij}\quad(i\ne j),\qquad B_{ii}=0.
\label{eq:kron-map}
\end{align}
The physical-bus input map is $\mathbf a(b)=\mathbf A_{\mathrm{inj}}\mathbf e_b$
in the stated ordering. With no algebraic buses, use $\mathbf L_{\mathrm{red}}=\mathbf L$
and $\mathbf A_{\mathrm{inj}}=\mathbf I_m$. The linear Kron map combined with
sinusoidal coupling is a reduced-model approximation. Network data follow
the IEEE test cases in MATPOWER~\cite{zimmerman2011matpower}.




\subsection{Active IEEE9 Configuration}
The active non-reset backend is IEEE9: nine electrical buses, three retained
dynamic machines and six uncertain $M_i/K_i$ coordinates. The new experiment
uses fixed physical probe and measurement bus 1, probe amplitude $0.05$~pu,
continuous durations $[0.2,3]$~s, separation $0.01$~s and recording window
$W=4$~s. With $T=3$, probing lasts 12~s. The five frequency samples occur
at $\{0.0015625,1.0015625,2.0015625,3,4\}$~s within each window and use
independent noise with $\sigma=0.005$~Hz.

The control magnitude is continuous on $[0,0.5]$~pu. The configured contingency
is a $-0.45$~pu step at physical bus 1, beginning immediately after probing;
control is injected at the same bus for the following 10~s. The monitored
limits are $\underline f=-0.2$~Hz and $\overline R=22$~Hz/s. These values,
the simplified disturbance and coverage $q=0.9$ are benchmark assumptions,
not cited operating standards or engineering acceptance certificates.
The ODE step is $1/640$~s; the control root solver uses 33 scan points and
$10^{-5}$~pu bracket tolerance. ODE discretization and root-bracket error
are distinct from posterior approximation error.


\begin{table}[t]
\caption{IEEE9 independent uniform prior bounds and known damping. Values are rounded here; resolved YAML values are retained with each run. Units follow Eq.~\eqref{eq:swing-frequency}.}
\centering
\begin{tabular}{lrrr}
\toprule Parameter & Machine 1 & Machine 2 & Machine 3\\
\midrule
$M_i^-$ & 0.0877899 & 0.0237671 & 0.0111780\\
$M_i^+$ & 0.1630383 & 0.0441390 & 0.0207591\\
$K_i^-$ & 0.029375 & 0.013125 & 0.007500\\
$K_i^+$ & 0.293750 & 0.131250 & 0.075000\\
$D_i$ & 0.0125414 & 0.0033953 & 0.0015969\\
\bottomrule
\end{tabular}
\end{table}
The prior is uniform on the product set in Eq.~\eqref{eq:latent-space};
its coordinate widths are study choices, not identification results.
The benchmark begins with zero relative rotor angles and zero speed deviations,
with nominal incremental power balance consistent with that equilibrium.
This operating point is a reduced perturbation-model assumption, not a
validated loaded AC dynamic operating point.

The full resolved priors, coupling, damping, simulator settings and seeds
must accompany every reported run. Changes to coverage or physical thresholds
must be declared before final testing and must not be selected to favor a method.
The obsolete six-duration catalogs and equilibrium master banks are not part
of this experiment.

\subsection{Verification and Comparison Scope}
A workability test uses $T=3$, training seed 101, evaluation seed 1001,
two training updates, training batch size four, four validation systems, two test systems,
eight posterior/planning particles and all six methods. It checks execution,
state carry, objective consistency and output completeness; its scores are
not performance evidence. Independent CPU/GPU trajectory checks, continuous
root-solver checks, posterior-truth exclusion and RL reward telescoping are
separate implementation checks.

A proposed full comparison uses independently trained policies for seeds
101/202/303, evaluated on seeds 1001--1005, with fixed declared optimization
budgets at each horizon. That plan is not a claim that the runs are complete.
Inference-particle sensitivity, EIG contrast sensitivity, numerical-gradient
sensitivity, convergence and checkpoint quality must be assessed before
publication-level ranking. The no-reset model and common budgets alone do
not prove that any method has converged or that a learned family wins.

Each metric gets a separate table with methods as rows and horizons as columns.
Use run-level sample standard deviations across the declared seeds;
per-system Bernoulli dispersion is not the uncertainty of a reported safety
rate. A single evaluation seed has no across-seed standard deviation.
Keep physical systems and their repeated noise outcomes together in paired
uncertainty calculations. Fixed training/calibration copies may be deduplicated
only if the resulting policies and evaluations are actually identical.
Report posterior MSC or MOCU as the primary objective-specific score, with
control, physical safety and the history-dependent oracle as appropriate.
Do not relabel realized regret as posterior MOCU. Timing must distinguish
offline optimization, online probe choice/refinement, terminal posterior/control
computation and evaluation-only oracle work.

\subsection{Other Systems and Historical Results}
SIR remains a separate EIG implementation benchmark. Its former results do
not validate the changed power-grid observation or control protocol.
IEEE14 requires a separately validated online state-carry and control backend
before comparison; the active implementation rejects unsupported execution.
The planned IEEE30 extension retains Demetriou et al.'s dynamic reference
with six machines, including two governed generators and four synchronous
condensers, and six added machine terminal buses~\cite{demetriou2017dynamic}.
Its proposed uncertainty is six inertia constants and two generator governor
parameters, rather than assigning $M/K$ to every electrical bus. The full
machine/exciter/governor implementation and initialization remain pending.
No non-reset IEEE14/IEEE30 results are claimed.

\section{Results and Remaining Validation}
\todo{Insert completed, matched non-reset multi-seed results separately for EIG, MSC and MOCU.}
Historical reset results and the earlier repeated-duration EIG pilot are
separate protocols. They must not populate the new comparison tables. A
successful smoke test establishes workability only, and no method superiority
is concluded from it.

\section{Discussion}
The probe sequence now changes both what is known about the grid and the
state from which terminal control acts. This is why the physical oracle is
history dependent and why a state-independent control bank is insufficient.
MSC directly values the selected control under a posterior safety constraint.
MOCU evaluates excess loss under its declared surrogate; EIG values parameter
information. These are different objectives, even where MSC and MOCU select
the same controller for a given history.

Engineering interpretation still requires validation of the reduced dynamics,
parameter priors, contingency, sensing assumptions, control authority and
physical limits. Neither coverage nor numerical agreement with a reduced
simulator certifies field safety. Zero decision latency, a single prescribed
contingency and a scalar step actuator limit the present application.
The MOCU shortfall loss is not a validated outage-damage model. Method claims
must also distinguish original algorithms from the documented objective and
numerical-optimization adaptations. Adaptivity and non-myopic gains remain
empirical questions rather than required outcomes imposed on the experiment.

\section{Conclusion}
The formulation connects an ordered sequence of noisy, non-reset probe
experiments to a continuous terminal power-injection decision. It separates
EIG, minimum safe control and finite-loss MOCU, and specifies the corresponding
physical safety evaluation. \todo{Add conclusions supported by completed
non-reset experiments after the required validation.}

\bibliographystyle{IEEEtran}
\bibliography{objective_driven_boed_refs}
\end{document}
